\documentclass[12pt]{article}
\usepackage[utf8]{inputenc}
\usepackage[T1]{fontenc}
\usepackage{hyperref}
\title{Sengseng Language Resources for K-12 Curriculum Development}
\author{Compiled by Claude Sonnet 5.5 (claude-sonnet-5-5) for Andr\'as Kornai}
\date{}
\begin{document}
\maketitle

\section{Introduction}
Sengseng (ISO 639-3: ssz) is an Austronesian language of the Pasismanua grouping, spoken in the south-west interior of West New Britain Province, Papua New Guinea. Sources give about 1,500 to 1,800 speakers, mostly from data around 2003. It is rated vulnerable or threatened. The web search found only general and encyclopedic material. It found no primer, textbook, dictionary or school curriculum material for Sengseng, and this report does not claim any exists.

\section{Primary Educational Resources}
No primers, textbooks or school materials for Sengseng turned up in the search. Sengseng is not known to be taught in schools. A curriculum project would have to build these materials from scratch, working with the community.

\section{Linguistic and Reference Materials}
\begin{itemize}
\item Wikipedia overview of the Sengseng language: \url{https://en.wikipedia.org/wiki/Sengseng_language}
\item Endangered Languages Project page, which has a bibliography section (I did not open it): \url{https://endangeredlanguages.com/lang/2650}
\item Max Planck Institute numeral study by Steve Henley (2006). It notes that Sengseng has a quinary counting system: \url{https://lingweb.eva.mpg.de/channumerals/Sengseng.htm}
\item Verbix page with a sample of Sengseng verb vocabulary: \url{https://docs.verbix.com/Languages/Sengseng}
\end{itemize}

\section{Religious and Community Resources}
\begin{itemize}
\item The Joshua Project profile describes the Sengseng people and their homeland near the Andru River: \url{https://joshuaproject.net/people_groups/14855}
\item Global Recordings Network language page for ssz. The search summary says it lists no audio recordings of its own. It reports Christian videos, Bibles and songs from another source, but I could not confirm the details: \url{https://www.globalrecordings.net/en/language/ssz}
\end{itemize}

\section{Digital and Audio Resources}
No Sengseng audio corpus, software or digital learning material was found in the search. The Global Recordings Network page above is the only audio-related lead, and it needs checking.

\section{Curriculum Development Context}
Sengseng is spoken by all adults but not all young people, and it is not known to be taught in schools. The documentation found is thin. A K-12 curriculum would need new fieldwork: a practical orthography, a word list, and recordings. The Endangered Languages Project bibliography is the first place to look for existing academic work. Related Pasismanua languages (Kaulong, Miu, Karore, Psohoh) might offer comparative material, but no resources for them were checked.

\end{document}